\phantomsection
\addcontentsline{toc}{subsection}{Restricción final}
\subsection*{Restricción final}

% Seccion con nombres de clase largos, que no se parten con guion.
\sloppy

\textit{Las operaciones deben poder ejecutarse en paralelo sobre diferentes cuentas (considera
concurrencia y estado mutable del Visitor). Además, el sistema debe garantizar que las
transacciones sobre una misma cuenta sean atómicas y consistentes (¿cómo lo resuelves con el
patrón elegido?).}

\vspace{0.4cm}

\subsubsection*{Qué pide, y por qué las dos mitades tiran en sentidos opuestos}

La restricción impone dos obligaciones que conviene leer juntas:

\begin{enumerate}[leftmargin=1.2cm]
    \item \textbf{Paralelismo} entre operaciones sobre cuentas \emph{distintas}.
    \item \textbf{Atomicidad y consistencia} entre operaciones sobre la \emph{misma} cuenta.
\end{enumerate}

Tomadas por separado, cada una tiene una solución trivial y equivocada: para la primera, no
sincronizar nada; para la segunda, un cerrojo global. Cada solución trivial destruye la otra
obligación. Lo que resuelve la tensión no es elegir, sino acertar con la
\textbf{granularidad}: el enunciado nombra la unidad de consistencia cuando dice ``sobre una
misma cuenta''. Si el cerrojo es por cuenta, las dos obligaciones se cumplen a la vez y sin
contradicción, porque dos operaciones sobre cuentas distintas no comparten cerrojo.

De ahí sale el criterio que guía toda esta sección: \textbf{la exclusión mutua pertenece al
elemento, no a la operación}. La cuenta es quien debe proteger sus invariantes; el visitante
solo pasa por ahí.

\subsubsection*{Tres intentos ingenuos, y qué rompe cada uno}

Antes de la solución conviene descartar lo que no funciona, porque las tres son tentaciones
razonables:

\begin{table}[h!]
    \centering
    \small
    \begin{tabular}{p{4.3cm}p{4.2cm}p{4.9cm}}
        \toprule
        \textbf{Intento} & \textbf{Qué arregla} & \textbf{Qué rompe} \\
        \midrule
        Un cerrojo global, o \texttt{lock} sobre el visitante completo
        & La consistencia: nunca hay dos operaciones a la vez.
        & El paralelismo exigido. Dos cuentas sin relación se serializan entre sí. \\
        \addlinespace
        Compartir una instancia de visitante entre hilos y acumular en sus campos
        & Nada. Parece ahorrar objetos.
        & El resultado: dos recorridos simultáneos mezclan sus acumuladores. \\
        \addlinespace
        Visitante sin estado, pero sin cerrojo en la cuenta
        & El estado del visitante, que deja de ser compartido.
        & Los invariantes de la cuenta: leer, sumar y escribir el saldo deja de ser atómico. \\
        \bottomrule
    \end{tabular}
\end{table}

El tercero es el más instructivo, porque es el que parece correcto. Un visitante puro no
arregla la cuenta: la carrera no está en el visitante sino en la secuencia
lectura--modificación--escritura del saldo, y esa vive en el elemento. La prueba 7 de la
ejecución lo mide.

\subsubsection*{Primera mitad: el estado mutable del visitante}

El patrón elegido tiene aquí un riesgo propio, y el enunciado lo señala al pedir que se
considere ``el estado mutable del Visitor''. Un visitante transporta datos, y la tentación es
que los acumule en sus campos. La implementación lo evita con tres decisiones, de la más
preferible a la tolerable:

\begin{itemize}[leftmargin=1.2cm]
    \item \textbf{Devolver en lugar de acumular.} La interfaz es genérica en el tipo de
          retorno, \texttt{IVisitanteOperacion<R>}, precisamente para que cada método de
          visita \emph{devuelva} un valor inmutable.
          \texttt{AuditarTransaccionVisitor} no tiene ni un campo: construye y devuelve un
          \texttt{InformeAuditoria}. Por eso una sola instancia se comparte entre hilos sin
          sincronización de ningún tipo.
    \item \textbf{Estado inmutable y de un solo uso.}
          \texttt{EjecutarTransaccionVisitor} sí lleva estado, la \texttt{Transaccion}, pero
          es un \textit{record} inmutable y el visitante se construye uno por operación. El
          estado existe, pero nadie lo escribe después de construirlo.
    \item \textbf{Si la acumulación compartida fuera inevitable}, contadores atómicos
          (\texttt{Interlocked}) o colecciones concurrentes. Nunca sincronizar el visitante
          completo: eso serializa los recorridos y devuelve al primer intento ingenuo.
\end{itemize}

Esta decisión es la que paga el inciso (c): un visitante que devuelve su resultado se puede
reducir, y un recorrido reducible se puede repartir entre hilos.

\subsubsection*{Segunda mitad: atomicidad por cuenta con el patrón elegido}

La pregunta del enunciado es concreta: \emph{¿cómo lo resuelves con el patrón elegido?} La
respuesta tiene tres piezas.

\textbf{1. El cerrojo vive en la cuenta.} Cada \texttt{Cuenta} lleva su propio objeto de
cerrojo. No es un detalle de implementación: es la afirmación de que la cuenta es la unidad de
consistencia del dominio.

\textbf{2. Un decorador lo toma y lo libera.} Para que ningún visitante concreto sepa que
existen cerrojos, el control de concurrencia entra como un \textbf{Decorator} del visitante,
no como código dentro de las operaciones:

{\small
\begin{verbatim}
public sealed class VisitanteTransaccional<R>
    : IVisitanteOperacion<R>
{
    private readonly IVisitanteOperacion<R> _interno;

    public R VisitAhorro(CuentaAhorro c)
        => BajoCerrojo(c, () => _interno.VisitAhorro(c));

    private static R BajoCerrojo(Cuenta c, Func<R> accion)
    {
        lock (c.Cerrojo)       // Monitor.Enter/Exit con
        {                      // try/finally, y reentrante
            return accion();
        }
    }
}
\end{verbatim}
}

Aquí se ve por qué el patrón elegido ayuda en lugar de estorbar: como la operación es un
objeto, se puede \emph{envolver}. Con las operaciones escritas como métodos de las cuentas, el
control transaccional habría tenido que repetirse dentro de cada método de cada cuenta, o
subir al cliente, que es donde peor está. Esta es una ventaja de Visitor que la métrica del
inciso (d) no mide.

\textbf{3. Orden total de adquisición cuando hacen falta dos cerrojos.} Una transferencia toca
dos cuentas. Dos transferencias cruzadas que tomaran los cerrojos en el orden de sus
argumentos podrían esperarse mutuamente para siempre. Se toman siempre por identificador
ascendente:

{\small
\begin{verbatim}
Cuenta primera = string.CompareOrdinal(origen.Id, destino.Id) <= 0
    ? origen : destino;
Cuenta segunda = ReferenceEquals(primera, origen) ? destino : origen;

lock (primera.Cerrojo)
{
    lock (segunda.Cerrojo)
    {
        ...
\end{verbatim}
}

El argumento de corrección es corto: un interbloqueo exige un ciclo en el grafo de espera, y
un ciclo exige que dos hilos adquieran dos cerrojos en órdenes opuestos. Con un orden total
sobre los identificadores, ese orden opuesto no puede ocurrir. No hace falta detección de
ciclos ni tiempos límite.

\subsubsection*{Hasta dónde llega la garantía}

Conviene decir con precisión qué queda garantizado, porque es menos de lo que la palabra
``atómico'' sugiere:

\begin{itemize}[leftmargin=1.2cm]
    \item \textbf{Sí:} toda operación sobre una cuenta es atómica frente a otras operaciones
          sobre \emph{esa} cuenta, y consistente, en el sentido de que ningún entrelazado
          puede dejar el saldo fuera de los límites que la cuenta impone.
    \item \textbf{Sí:} operaciones sobre cuentas distintas progresan en paralelo, sin
          contención entre ellas.
    \item \textbf{No:} una operación sobre un \emph{grupo} de cuentas no es atómica. El
          decorador da atomicidad por cuenta, y el inciso (c) lo muestra con una aplicación
          parcial real. Tomar todos los cerrojos del grupo da \textbf{aislamiento}, no
          atomicidad: si la operación falla a mitad, lo ya aplicado sigue aplicado. Para
          todo\,o\,nada hace falta una validación previa o compensación por \textit{saga}.
    \item \textbf{No:} no hay durabilidad. Nada se escribe en un registro antes de aplicarse,
          así que un fallo del proceso pierde el estado. Un sistema real delegaría esto en las
          transacciones de la base de datos o en un registro de escritura anticipada.
    \item \textbf{No:} \texttt{lock} no es justo. No garantiza orden de llegada entre hilos
          que esperan. Si el dominio exigiera orden FIFO por cuenta, la forma correcta sería
          una cola por cuenta, es decir un actor por cuenta, que convierte la serialización en
          una propiedad estructural en lugar de depender de cerrojos.
\end{itemize}

\subsubsection*{Evidencia de ejecución}

Las partes 7 y 8 de \texttt{Demo} miden exactamente las dos mitades de la restricción. Ocho
hilos hacen 250 depósitos de 10,00 sobre la misma cuenta, y después dieciséis hilos operan
sobre cuatro cuentas a la vez:

{\small
\begin{verbatim}
7. Restriccion: concurrencia sobre la MISMA cuenta
--------------------------------------------------
  con decorador transaccional: esperado 20,000.00, obtenido 20,000.00
  atomico y consistente: true
  sin decorador:               esperado 20,000.00, obtenido 3,390.00
  actualizaciones perdidas:    16,610.00

8. Restriccion: paralelismo sobre cuentas DISTINTAS
---------------------------------------------------
  PAR-AHO     saldo=  10,000.00 correcto=true
  PAR-COR     saldo=  10,000.00 correcto=true
  PAR-INV     saldo=  10,000.00 correcto=true
  PAR-EMP     saldo=  10,000.00 correcto=true
  16 hilos sobre 4 cuentas, sin contencion entre ellas: 5 ms
\end{verbatim}
}

Dos lecturas de estas cifras. La primera: con el decorador el saldo final es exacto, y sin él
se pierde más de la mitad del dinero depositado. La segunda, menos cómoda: \textbf{la cifra de
pérdida cambia en cada ejecución}, porque depende del entrelazado de los hilos. Es un error no
determinista, y eso es justamente lo que lo hace peligroso: una prueba que corriera una sola
vez podría pasar con el código incorrecto. Para que la carrera se manifieste en cada corrida,
\texttt{Cuenta} amplía a propósito la ventana entre la lectura y la escritura del saldo con un
\texttt{Thread.Yield()} activado solo en esa prueba; en producción no existiría.

\subsubsection*{Resumen del proceso}

El camino seguido para resolver el ejercicio, que es el mismo que serviría para cualquier otro
problema de diseño:

\begin{enumerate}[leftmargin=1.2cm]
    \item Leer el enunciado buscando los \emph{ejes de variación}, no las clases: aquí, tipos
          de cuenta y operaciones.
    \item Identificar los conceptos del dominio y cuál de los ejes está cerrado.
    \item Escribir mentalmente la solución ingenua, con \textit{switch} sobre el tipo de
          cuenta, y nombrar qué rompe.
    \item Buscar el patrón que cierra el eje correcto, descartando los candidatos con
          argumentos y no por eliminación (inciso a).
    \item Diseñar las dos jerarquías y el mecanismo de doble despacho (inciso a).
    \item Implementar las $T \times V$ combinaciones de modo que cada celda sea una regla de
          negocio real (inciso b).
    \item Integrar con los patrones estructurales que el problema pide, pagando y nombrando su
          peaje (inciso c).
    \item Medir la decisión sobre el propio repositorio en lugar de afirmarla (inciso d).
    \item Resolver la restricción eligiendo la granularidad correcta, y decir hasta dónde llega
          la garantía (esta sección).
\end{enumerate}

\fussy
